% AulaTeX - plantilla materializada desde memoria editorial
\documentclass[12pt,letterpaper,oneside]{article}
\usepackage[margin=2.35cm]{geometry}
\usepackage[utf8]{inputenc}
\usepackage[T1]{fontenc}
\usepackage[spanish,es-tabla]{babel}
\usepackage[scaled=.96]{helvet}
\renewcommand{\familydefault}{\sfdefault}
\usepackage{setspace}
\usepackage{graphicx}
\usepackage[table]{xcolor}
\usepackage{booktabs}
\usepackage{array}
\usepackage{longtable}
\usepackage{hyperref}
\usepackage[authoryear,round]{natbib}

\definecolor{unadmGreenDark}{HTML}{174A3A}
\definecolor{unadmGreen}{HTML}{5F8F3A}
\definecolor{unadmGold}{HTML}{B88A2A}
\definecolor{unadmPaper}{HTML}{F6F7F2}
\definecolor{unadmInk}{HTML}{1F2A24}
\hypersetup{colorlinks=true,linkcolor=unadmGreenDark,citecolor=unadmGreenDark,urlcolor=unadmGreenDark}
\setlength{\parindent}{0pt}
\setlength{\parskip}{0.72em}
\onehalfspacing
\bibliographystyle{plainnat}

\newcommand{\pendiente}[1]{\textcolor{red}{[PENDIENTE: #1]}}

\begin{document}

\begin{titlepage}
  \pagecolor{unadmPaper}
  \color{unadmInk}
  \begin{center}
    \vspace*{0.8cm}
    \includegraphics[height=2.05cm]{img/departamentos/UnADM.pdf}\\[0.7cm]
    {\Large Universidad Abierta y a Distancia de Mexico}\\[0.2cm]
    {\large Licenciatura en Derecho}\\[1.2cm]
    {\color{unadmGreenDark}\rule{0.82\linewidth}{1.4pt}}\\[0.55cm]
    {\Huge\bfseries Actividad 10 - Seminario I\par}
    \vspace{0.35cm}
    {\Large Actividad 10 - Seminario I\par}
    \vspace{0.45cm}
    {\color{unadmGold}\rule{0.62\linewidth}{1.1pt}}\\[1.1cm]
    \begin{tabular}{rl}
      \textbf{Alumno:} & Martin Jonathan de la Cruz \\
      \textbf{Matricula:} & ES2611202040 \\
      \textbf{Figura docente:} & Nombre por definir \\
      \textbf{Materia:} & Seminario I \\
      \textbf{Semestre/Bloque:} & 2 / 1 \\
      \textbf{Tipo/Creditos:} & Obligatoria / 8 \\
      \textbf{Fecha:} & \today
    \end{tabular}
    \vfill
    {\small Roma Norte, Ciudad de Mexico}
  \end{center}
  \nopagecolor
\end{titlepage}

\tableofcontents
\clearpage

\section*{Resumen editorial}
\addcontentsline{toc}{section}{Resumen editorial}
Esta plantilla materializa la memoria editorial de \textit{Seminario I}.
La entrega debe convertir la consigna en un problema juridico delimitado, con
fundamento verificable, analisis propio y conclusion aplicable.

\section{Introduccion}
\pendiente{Abrir con una afirmacion fuerte que traduzca la consigna a problema juridico.}

\section{Encuadre institucional y juridico}
La materia se trabaja desde la identidad academica de la UnADM y desde la
formacion juridica aplicada. La bibliografia local debe sostener el analisis
con fuentes institucionales, normativas y academicas verificables \citep{unadmSitioWeb,unadmMallaDerecho2024}.

\section{Memoria editorial aplicada}
\begin{itemize}
  \item Organizar la entrega en problema, marco juridico, analisis, producto, conclusion y referencias.
\end{itemize}

\section{Marcadores de investigacion}
\begin{itemize}
  \item Delimitar fuentes oficiales, doctrina pertinente y criterios aplicables antes de redactar.
\end{itemize}

\section{Desarrollo del producto}
\pendiente{Insertar aqui el producto solicitado: informe, cuadro, mapa, analisis de caso, matriz, linea de tiempo o ensayo.}

\subsection{Analisis propio}
\pendiente{Explicar la consecuencia juridica del producto y sostener una postura personal argumentada.}

\section{Checklist de calidad}
\begin{itemize}
  \item Verificar citas, bibliografia, coherencia de estructura, ortografia y compilacion.
\end{itemize}

\section{Conclusion}
\pendiente{Cerrar con aprendizaje, criterio juridico propio y aplicacion profesional.}

\clearpage
\nocite{unadmSitioWeb,unadmMallaDerecho2024}
\bibliography{seminario-i}

\end{document}
